% Generated draft: requires mathematical and source review.
\documentclass[11pt]{article}
\usepackage[margin=25mm]{geometry}
\usepackage{fontspec}
\setmainfont{Latin Modern Roman}
\usepackage{amsmath,amssymb,mathtools,unicode-math}
\setmathfont{Latin Modern Math}
\usepackage{enumitem,needspace,xurl,hyperref}
\setlength{\parindent}{0pt}
\setlength{\parskip}{5pt}
\setlength{\emergencystretch}{3em}
\newcommand{\N}{\mathbb{N}}
\newcommand{\Z}{\mathbb{Z}}
\newcommand{\Q}{\mathbb{Q}}
\newcommand{\R}{\mathbb{R}}
\newcommand{\C}{\mathbb{C}}
\newcommand{\F}{\mathbb{F}}
\newcommand{\A}{\mathbb{A}}
\newcommand{\PP}{\mathbb{P}}
\newcommand{\E}{\mathbb{E}}
\newcommand{\eps}{\varepsilon}
\newcommand{\dd}{\,\mathrm{d}}
\begin{document}
\section*{UnsolvedMath ID: 700035}
\textbf{Source problem code:} catalogue entry
\subsection{Definitions and mathematical statement}
For a group \(H\), use the commutator convention \([x,y]=x^{-1}y^{-1}xy\), and define iterated commutators by
\[
[x,{}_1y]=[x,y],\qquad [x,{}_{n+1}y]=[[x,{}_ny],y]\quad(n\geq 1).
\]
An element \(y\in H\) is an Engel element of \(H\) if, for every \(x\in H\), there exists an integer \(n=n(x,y)\geq 1\) such that \([x,{}_ny]=1_H\). A \emph{nilgroup}, in the terminology of the supplied source, is a group all of whose elements are Engel elements. Thus \(H\) is a nilgroup precisely when
\[
\forall x,y\in H\;\exists n\geq 1:\quad [x,{}_ny]=1_H.
\]
Here a nilsubgroup is understood to be a subgroup that is itself a nilgroup, with the Engel condition evaluated within that subgroup.

The question is whether the following implication holds for every group \(G\) and every pair of normal subgroups \(A,B\trianglelefteq G\):
\[
\left.
\begin{gathered}
G=AB=\{ab:a\in A,\ b\in B\},\\
\forall H\in\{A,B\}\;\forall x,y\in H\;\exists n\geq 1:\ [x,{}_ny]=1_G
\end{gathered}
\right\}
\quad\Longrightarrow\quad
\forall x,y\in G\;\exists n\geq 1:\ [x,{}_ny]=1_G.
\]
No finiteness assumption or uniform bound on the Engel lengths is imposed. The source does not separately define ``nilsubgroup''; the intrinsic subgroup interpretation used above should be checked against the original terminology.
\subsection{Short English statement}
Must every group that is the product of two normal subgroups, each itself an Engel group, also be an Engel group? The required Engel length may depend on both elements. This uses the intrinsic interpretation of the source's term ``nilsubgroup.''
\subsection{Sources}
\begin{itemize}
\item Local UnsolvedMath registry: supplied complete record, permanent ID \(700035\), exact source code \texttt{catalogue entry}, title \emph{Kourovka Notebook Problem 1.40}; upstream revision \texttt{d371ee917ac9647fbb06cd6fc529deec6f3dda82}.
\item Supplied attribution: Sh. S. Kemkhadze, proposed in 1965; original problem number 1.40 in \emph{The Kourovka Notebook (Unsolved Problems in Group Theory)}, first issue (1965). The record identifies the consulted edition as the 21st issue edition (2026), \texttt{arXiv:1401.0300v48}. The collection description names E. I. Khukhro and V. D. Mazurov as editors.
\item The registry labels the problem \texttt{open} and records an update at \texttt{2026-10-16T14:00:00Z}. These are supplied registry metadata, not an independently verified literature assessment. The referenced edition has not been independently consulted for this draft.
\end{itemize}
\end{document}
